\scanpage{31}
\mspage{5}
canonique stable par translations etc, mais une classe de telles mod. homothétie par $c\in\mathbb R^{*+}$ - ces métriques sont caractérisées par la condition d'être invariantes par translation [\uline{NB} $X$ a automatiquement une structure de droite affine sur $\mathbb C$]. Ces métriques \deleted{\unsure{sont définies}} ne dépendent que de la structure conforme sous-jacente.

3) \uline{Cas} \deleted{\unsure{de la sphère} $X\simeq\mathbb P^1_{\mathbb C}$} $X\simeq$ \unsure{plan} projectif réel. Dans ce \deleted{premier} \deleted{\unsure{cas}} cas, $X^{or}$ est une droite projective complexe munie d'un anti-automorphisme invol. $\sigma$ \uline{sans pt fixe} (antipodisme), ou ce qui revient au $=$, d'une métrique ``euclidienne sphérique'' compatible avec sa structure conforme. [Celle-ci, \deleted{\unsure{étant stable}}\inserted{\unsure{se normalise}} \deleted{par l'antipodisme,}\inserted{par la condition que l'aire de la surface soit $=1$}] Étant invariante par l'antipodisme, cette métrique se descend à $X$.

$3'$) Cas $X\simeq\mathbb P^1_{\mathbb C}$. Il n'y a pas de métrique riemannienne sur $X$, invariante [ou quasi-invariante] par tt automorphismes - \unsure{donc sans} structure suppl. sur $X$
% archive foliation: 33
